\documentclass[12pt]{article}
\usepackage[utf8]{inputenc}
\usepackage[T1]{fontenc}
\usepackage{hyperref}
\title{Naukan Yupik Language Resources for K-12 Curriculum Development}
\author{Compiled by Claude Sonnet 5.5 (claude-sonnet-5-5) for Andr\'as Kornai}
\date{}
\begin{document}
\maketitle

\section{Introduction}
Naukan Yupik (ISO 639-3 ynk) is an Eskimo-Aleut language of the Russian Far East, historically spoken at Naukan (Nuvuqaq) on the Bering Strait. The village was closed by the Soviet authorities in 1958. A review in \textit{\'Etudes/Inuit/Studies} reports that only about 60 middle-aged and elderly people, out of a Nuvuqaqmiit population of about 450, are still fluent. The language is not known to be taught in schools. The web searches found very little beyond the dictionary described below, and no primer, textbook or grammar.

\section{Primary Educational Resources}
No primers, textbooks or school materials for Naukan were found in the search results.

\section{Linguistic and Reference Materials}
\begin{itemize}
\item \textit{Naukan Yupik Eskimo Dictionary} (2004), compiled by Dobrieva, Golovko, Jacobson and Krauss, edited by Jacobson, and published by the Alaska Native Language Center. It has two parallel volumes: a Latin-orthography to English volume, and a volume in the modified Cyrillic alphabet used for the Chukotkan Eskimo languages with Russian glosses. Review and description by Jacobson in \textit{\'Etudes/Inuit/Studies} 2005: \url{https://id.erudit.org/iderudit/013937ar} and \url{https://www.erudit.org/en/journals/etudinuit/2005-v29-n1-2-etudinuit1430/013937ar/abstract}
\item Related article in the same issue: \url{https://www.erudit.org/fr/revues/etudinuit/2005-v29-n1-2-etudinuit1430/013953ar.pdf}
\item A Bering Straits language-resources page titled ``yupik eskimo dictionary vol 2''. The search result did not make clear which language it covers, so verify before use: \url{https://beringstraits.com/language-resources/yupik_eskimo_dictonary_vol_2}
\item The WALS Online record for Yupik (Naukan): \url{https://openpolar.no/Record/ftolac:oai:wals.info:yna}
\end{itemize}

\section{Religious and Community Resources}
None were found in the search results.

\section{Digital and Audio Resources}
No audio or digital learning resources for Naukan were found in the search results. A blog overview of the languages of the Bering Strait region is at \url{https://thelanguagecloset.com/2024/08/31/the-languages-of-asias-easternmost-frontier-pt-2/}

\section{Curriculum Development Context}
With about 60 elderly fluent speakers and no known school teaching, any curriculum would have to be built from the 2004 dictionary and from documentation work with the remaining speakers. The dictionary gives both a Latin and a Cyrillic orthography, so orthography choice is a first decision. The related Central Siberian Yupik and Central Alaskan Yup'ik may offer teaching models, but they are different languages and were not researched here.
\end{document}